\documentclass[10pt,a4paper]{amsart}
\usepackage[margin=19mm]{geometry}
\usepackage{amsmath,amssymb,mathtools}
\usepackage[scr=rsfs,cal=euler]{mathalfa}
\usepackage{xfrac}
\usepackage[unicode,hidelinks,bookmarks=false]{hyperref}
\newcommand{\ie}{i.e.\ }
\newcommand{\cf}{cf.\ }
\newcommand{\resp}{respectively\ }
\newcommand{\Subsection}{\subsection}
\providecommand{\nobreakdash}{\nobreak\hbox{-}}
\setlength{\emergencystretch}{2em}
\multlinegap0pt
\usepackage{enumerate}
\usepackage[shortlabels]{enumitem}
\newtheorem{lem}{Lemma}
\newtheorem{claim}{Claim}
\newcommand{\rmv}{\smallsetminus}
\newcommand{\set}[1]{ \left\{ #1 \right\} }
\title{Genus two Goeritz equivalence in lens spaces $L(p, 1)$}
\author{Brandy Doleshal, Matt Rathbun}
\date{Source: arXiv:2602.16458v2 (CC BY 4.0)}
\begin{document}
\maketitle

\noindent\emph{Source and licence.} Genus two Goeritz equivalence in lens spaces $L(p, 1)$. Version arXiv:2602.16458v2 (CC BY 4.0), distributed by arXiv under the Creative Commons Attribution 4.0 International licence, http://creativecommons.org/licenses/by/4.0/, as recorded in the arXiv licence metadata for this exact version and shown on the abstract page https://arxiv.org/abs/2602.16458v2. Full licence notice: \texttt{licenses/arxiv-2602.16458-CC-BY-4.0.txt}.\par\smallskip

\noindent\textit{Editorial scope.} Counted unit: the article's lem lem:inequalities (source0.tex lines 626--627). The article's complete original proof of it is delivered with it (source0.tex lines 629--678, one \texttt{proof} environment of 50 lines). No mathematics is written, repaired or re-derived: the statement and the proof are the article's own lines, in the article's own notation, and no retained line is substituted or re-flowed.  No further source-local context is needed: every cross-reference of every retained line resolves inside the two delivered spans.  Nothing was omitted from any retained span, so the package carries no omission record.  No retained line contains a citation, so the wrapper carries no bibliography and no bibliography entry.  No retained line contains a figure environment, an image-inclusion command or drawing code, so no image is delivered and the package declares no asset dependency.  Editorial formatting only, in addition: an AMS \texttt{amsart} preamble that restates the article's own declarations and macros used by the retained lines, namely the environments \texttt{lem}, \texttt{claim} on separate counters, together with the standard packages those lines need.  The retained environments print in this document as Lemma 1, Claim 1 in order of appearance, whereas the article prints the same items with the numbering of its own full document.  The complete native source of the article as distributed in the arXiv e-print for this exact version is bundled unchanged as \texttt{sources/194883/source0.tex}.
\begin{lem} \label{lem:inequalities} For all $p \geq 4$,  \[\mathcal{H}_p \subseteq  \mathcal{B}_p .\]
\end{lem} 

\begin{proof} First, observe that the map from $\mathcal{G}_p$ to $\mathcal{H}_p$ factors through \[\mathcal{G}_p'  = \mathcal{G}_p / \langle \langle \beta^2 \rangle \rangle. \]


In this case, $p \geq 4$, so $\mathcal{G}_p' = \langle \, \alpha \mid \alpha^2 = 1 \, \rangle \oplus \langle \, \beta, \, \gamma, \, \sigma \mid \beta^2 = \gamma^2 = \sigma^2= 1 \, \rangle$. We will abuse notation slightly and refer the generators in $\mathcal{G}_p'$ again as simply $\alpha$, $\beta$, $\gamma$, and $\sigma$.

To prove this lemma, we will induct on the reduced word length of elements $w'$ of $\mathcal{G}_p'$. Without loss of generality, we may ignore $\alpha$, and assume that all matrices are products of images of $\beta$, $\gamma$, and $\sigma$. Consider a matrix $A = \begin{pmatrix} a & b \\ c & d \end{pmatrix} \in \mathcal{H}_p$ as a M\"{o}bius transformation, acting on the extended real line, $\overline{\mathbb{R}} = \mathbb{R} \cup \set{\infty}$. Observe that $q(\gamma_*)(x) = -x - 1$, $q(\sigma_*)(x) = \frac{x}{-px - 1}$, and $q(\beta_*)(x) = -x$. For a reduced nontrivial word $w' \in \mathcal{G}_p'$, let $LL(w')$ refer to the smallest, left-most subword that contains either a $\gamma$ or a $\sigma$. Thus, $LL(w') \in \set{\gamma, \sigma, \beta \gamma, \beta \sigma, \emptyset }$, and $LL(\beta) = \emptyset$. 

Observe that $q(\gamma_*)$ has fixed points $-\frac{1}{2}$ and $\infty$, and exchanges the two complementary intervals of $\overline{\mathbb{R}} \rmv \set{-\frac{1}{2}, \infty}$. Similarly, $q(\sigma_*)$ has fixed points $-\frac{2}{p}$ and $0$,  $q(\beta_*)$ has fixed points $0$ and $\infty$, and both exchange the two complementary intervals of their respective fixed points. So, let 
\[ I_\gamma^- = \left( -\infty, -\frac{1}{2} \right] \cup \set{\infty}, \qquad I_\gamma^+ = \left[\frac{1}{2}, \infty \right) \cup \set{\infty}, \qquad I_\gamma  = I_\gamma^- \cup I_\gamma^+,  \]
%
and
%
\[ I_\sigma^- = \left[ -\frac{2}{p}, 0 \right], \qquad I_\sigma^+ = \left[ 0, \frac{2}{p} \right], \qquad I_\sigma = I_\sigma^- \cup I_\sigma^+.\]
%
Finally, set $P = \set{0, \infty}$.

We make an initial observation. If $p > 4$, then $I_\gamma$ and $I_\sigma$ are disjoint. If $p = 4$, then $\frac{2}{p} = \frac{1}{2}$, and the two intervals $I_\sigma$ and $I_\gamma$ overlap at endpoints. However,

\begin{claim}  \label{claim:overlap} If $p = 4$,  $q(w_*)(0) \neq \pm \frac{1}{2}$  and $q(w_*)(\infty) \neq \pm \frac{1}{2}$. 
\end{claim}

\begin{proof}
If $q(w_*) = A = \begin{pmatrix} a & b \\ c & d \end{pmatrix}$, then $A(0) = \frac{b}{d}$ and $A(\infty) = \frac{a}{c}$. But if $\frac{b}{d} = \pm \frac{1}{2}$, this would mean that $d = \pm 2b$, and if $\frac{a}{c} = \pm \frac{1}{2}$, then $c = \pm 2a$. However, since $\pm A \in \mathcal{A}_p$, $\pm a = kp + 1$ and $\pm b = \ell$ for some $k, \ell \in \mathbb{Z}$, $\pm c = mp$, and $\pm d = np + \Delta$ for some $m, n \in \mathbb{Z}$. Neither of these are then possible if $p = 4$.
\end{proof}

\begin{claim} \label{claim:images} For a reduced nontrivial word $w \in \mathcal{G}_p$, 
\begin{enumerate}
\item \label{case:G} if $LL(w') = \gamma$, then $q(w_*)(P) \subseteq I_\gamma^-$, and 
\item \label{case:S} if $LL(w') = \sigma$, then $q(w_*)(P) \subseteq I_\sigma^-$.
\end{enumerate}

\end{claim}

\begin{proof} We prove the claim by induction on reduced wordlength of words in $\mathcal{G}_p'$. 

Suppose $w \in \mathcal{G}_p$ and $|w'|= 1$. If $LL(w') = \gamma$, then, $q(w_*) = q(\gamma_*)$. Thus, $q(\gamma_*)$ fixes $\infty$, and $q(\gamma_*)(0) = -1 \in I_\gamma^-$, which establishes (\ref{case:G}). If $LL(w') = \sigma$, then $q(w_*) = q(\sigma_*)$. Thus, $q(\sigma_*)$ fixes $0$, and $q(\sigma_*)(\infty) = -\frac{1}{p} \in I_\sigma^-$, which establishes (\ref{case:S}).

If $w \in \mathcal{G}_p$ and $|w'| = 2$, then there are 6 possible words that $w'$ can represent. Four of them satisfy one of the hypotheses of Claim \ref{claim:images}: $w' = \gamma \sigma$ and $\gamma \beta$ satisfy the hypothesis of (\ref{case:G}), while  $w' = \sigma \gamma$ and $w' =  \sigma \beta $ satisfy the hypothesis of (\ref{case:S}). Each of these may be checked directly, as $q(\gamma_* \sigma_*) = \begin{pmatrix} -p+1 & -1 \\ p & 1 \end{pmatrix}$, $q(\gamma_* \beta_*)= \begin{pmatrix} 1 & -1 \\ 0 & 1 \end{pmatrix}$, $q(\sigma_* \gamma_*) = \begin{pmatrix} 1 & 1 \\ -p & -p+1 \end{pmatrix}$, and $q(\sigma_* \beta_*)= \begin{pmatrix} 1 & 0 \\ -p & 1 \end{pmatrix}$.

Now, assume that for $n > 1$, the claim holds for all words with length less than $n$, and suppose $w' \in \mathcal{G}_p'$ with $|w'| = n$. 
In order for the hypotheses of Claim \ref{claim:images} to apply, it must be that $w' = gv$, for some $v \in \mathcal{G}_p'$ with $|v| = n-1$, and $g \in \set{\gamma, \sigma}$. 

If $g = \gamma$, then $LL(v) \in \set{\sigma,  \beta \sigma,  \beta \gamma, \emptyset}$. Suppose $LL(v) = \emptyset$. Then, $v = \beta$, so $q(v_*)$ fixes both $0$ and $\infty$, so $q(w_*)(P) = q(\gamma_*)(P) \subseteq I_\gamma^-$. Suppose $LL(v) \in \set{\sigma, \beta \sigma}$. Then, inductively, $q(v_*)(P) \in I_\sigma^{\pm}$, and by Claim \ref{claim:overlap}, $q(v_*)(P)$ is in the complement of $[-\infty, -\frac{1}{2}]$, so that $q(w_*)(P) \subseteq I_\gamma^-$. Suppose finally that $LL(v) = \beta \gamma$. Then inductively, $q(v_*)(P) \subseteq I_\gamma^+$, so that $q(w_*)(P) \subseteq I_\gamma^-$. 

Similarly, if $g = \sigma$, then
$LL(v) \in \set{\gamma, \beta \gamma, \beta \sigma, \emptyset}$. Suppose $LL(v) = \emptyset$. Then, $v = \beta$, so $q(v_*)$ fixes both $0$ and $\infty$, so $q(w_*)(P) = q(\sigma_*)(P) \subseteq I_\sigma^-$. Suppose $LL(v) \in \set{\gamma, \beta \gamma}$. Then, inductively, $q(v_*)(P) \in I_\gamma^{\pm}$, and by Claim \ref{claim:overlap}, $q(v_*)(P)$ is in the complement of $[-\frac{2}{p}, 0]$, so that $q(w_*)(P) \subseteq I_\sigma^-$. Suppose, finally, that $LL(v) = \beta \sigma$. Then, inductively, $q(v_*)(P) \subset I_\sigma^+$, so that $q(w_*)(P) \subseteq I_\sigma^-$.
\end{proof}

We are now in a position to proceed with the proof of Lemma \ref{lem:inequalities}. Let $w \in \mathcal{G}_p$ with reduced word $w' \in \mathcal{G}_p'$, so that $q(w_*) = \begin{pmatrix} kp + 1 & \ell \\ mp & np + \Delta \end{pmatrix} \in \mathcal{H}_p$ (still ignoring $\alpha$). Observe that $q(w_*) \in \mathcal{B}_p$ if and only if $q(\beta_* w_*) \in \mathcal{B}_p$, so we may assume that $LL(w') \neq \beta$. If $LL(w') = \gamma$, then by Claim \ref{claim:images}, $q(w_*)(\infty) \in I_\gamma^-$. Observe that $q(w_*)(\infty) = \infty$ if and only if $q(w_*) = \begin{pmatrix} kp + 1 & \ell \\ 0 & np + \Delta \end{pmatrix}$, which is evidently in $\mathcal{B}_p$. On the other hand, if $q(w_*)(\infty) \neq \infty$, then since $q(w_*)(\infty) \in I_\gamma^-$, we have $|\frac{kp+1}{mp} | = |q(w_*)(\infty)| \geq \frac{1}{2}$, from which it follows that $|mp| \leq 2|kp+1|$. On the other hand, if $LL(w')  = \sigma$, then by Claim \ref{claim:images}, $q(w_*)(0) \in I_\sigma^-$. This means that $|\frac{\ell}{np+\Delta}| = |q(w_*)(0)| \leq \frac{2}{p}$, from which it follows that $p|\ell | \leq 2|np+\Delta|$. Thus, in all cases, $q(w_*) \in \mathcal{B}_p$.
\end{proof}

\end{document}
